\documentclass[11pt]{article}
\usepackage[T1]{fontenc}
\usepackage{amsmath,amssymb,amsthm}
\usepackage[margin=1in]{geometry}
\usepackage[hidelinks]{hyperref}
\newtheorem{theorem}{Theorem}
\newtheorem{lemma}[theorem]{Lemma}
\newtheorem{corollary}[theorem]{Corollary}
\newcommand{\DD}{d'd''}
\newcommand{\C}{\mathcal C}
\newcommand{\G}{\mathcal G}
\newcommand{\MA}{\operatorname{MA}}
\newcommand{\A}{\mathfrak A}
\title{An exact majorant criterion and a countability obstruction\\
for actual finite-energy presentations}
\author{Independent scoped P16.7 draft; no correctness grant}
\date{9 October 2026}
\begin{document}
\maketitle
\begin{abstract}
We give a necessary and sufficient recognition criterion in the entire
original continuous class. It separates a countable continuous-barrier
condition, expressed by cozero sets on the actual compact analytic space,
from a uniform lower energy bound for analytic majorants under every
strict continuous obstacle. Neither condition assigns positivity or an
operator to a raw limit. Their conjunction constructs an actual decreasing
presentation by an explicit summable-error selection. Every eligible
measure then has a finite moment against the recognized potential, with a
bound involving the same obstacle profile. This is a recognition bridge;
it does not establish the missing obstacle condition for the prescribed
isotropic Dirac, variational attainment, or the original equation theorem.
On the actual designated space we also prove that its unit Gauss point
is not a countable intersection of open sets. Its bounded upper
semicontinuous singleton indicator consequently admits no decreasing
continuous presentation at all, giving an obstruction for the entire
original generated class, without assigning it a psh status or an operator.
\end{abstract}

\section*{Exact datum and bounded structural input}
Retain the supplied compact, good, strict, smooth, separated, boundaryless
pure analytic surface
\[
 K=\mathbb C((s)),\quad |s|=e^{-1/8},\quad t=s^8,\qquad
 X=(\mathbb G_m^2)^{\rm an}/\langle(t,2),(3,t)\rangle.
\]
These geometric properties are inputs, not newly proved. Genuine
nonalgebraization and coherent descent remain uncertified. In local
coordinates $r_i=-\log|z_i|$, set $\rho=|r|^2/2$ and $\theta=\DD\rho$.
Deck changes translate $r$ by integers, so changes of $\rho$ are smooth
and $\DD$-null. The original raw mass $\int_X\theta^2=2$ is retained
as part of the datum, rather than derived from a new comparison theorem.

The ENTIRE original class $\C$ consists of real continuous weights $a$
whose local potentials $\rho+a$ belong to the minimal continuous CLD psh
sheaf: the smallest continuous sheaf containing smooth positive-Hessian
functions and closed under local uniform limits. No finite-log-max,
skeletal, radial, or classical substitute is made. Write
\[
 \theta_a=\theta+\DD a,\quad B(a,b)=\theta_a\wedge\theta_b,\quad
 E(a)=\frac13\int_X a\bigl(\theta^2+B(a,0)+B(a,a)\bigr).
\]
Here $\G$ means exactly functions admitting actual continuous sequences
$a_n\in\C$ decreasing everywhere to the function, with
$\inf_n E(a_n)>-\infty$. No presentation-generated operator theorem is
invoked in this paper, and no operator or energy is assigned to an
unrepresented function.

The bounded structural facts used below need no Green identity with two
nonsmooth factors and no energy concavity. We reconstruct them here.
The local continuous sheaf is stable under adding a constant: for fixed
$c$, the functions $f$ for which $f+c$ belongs to the sheaf form a
uniformly closed sheaf containing all smooth positive-Hessian generators.
Doing this for every $c$ proves $a+c\in\C$. Global uniform closure of
$\C$ follows immediately from its local definition.

CLD v2, Lemma 8.5.5 and construction 8.7.1 give positive closed
recursive products of continuous psh potentials. Section 8.8.5 gives
multilinearity for local differences of such functions and invariance
under smooth $\DD$-null changes. Thus $B(a,b)$ is globally defined.
Every weight is locally dpsh, since $a=(\rho+a)-\rho$ and $\rho$ is
smooth psh. The recursive derivative identity gives globally
\[
 B(a,b)-\theta^2=\DD(a\theta_b)+\DD(b\theta).
\]
Both sides glue as currents; the multipliers on the right are continuous
and bounded on compact $X$. Evaluate on the global compactly supported
smooth test $1$. A distributional derivative evaluates on $\DD1=0$;
the two exact terms therefore vanish. All bounded mixed masses are two.
Positive top currents are positive Radon measures by CLD 7.2.8, so the
displayed energy is finite for every $a\in\C$. A constant does not
change $\theta_a$. Consequently
\begin{equation}\label{shift}
                 E(a+c)=E(a)+2c.
\end{equation}
This uses only bounded recursive products and their defining derivative
identity, not an operator at a Borel function.

\section*{A countable barrier condition without metrizability}
A set $U\subset X$ is a \emph{cozero set} if
$U=\{f>0\}$ for a continuous function $f:X\to[0,1]$.

\begin{lemma}\label{barriers}
Let $u:X\to[-\infty,\infty)$ be bounded above. The following are
equivalent:
\begin{enumerate}
\item There exist real continuous $h_n\downarrow u$ everywhere.
\item For every rational $q$, the set $\{u<q\}$ is a cozero set.
\end{enumerate}
Either condition implies upper semicontinuity. No countable dense set,
metrizable carrier, or countable analytic test family is required.
\end{lemma}
\begin{proof}
If $h_n\downarrow u$, then $\{u<q\}=\bigcup_n\{h_n<q\}$.
The uniformly convergent continuous series
\[
 f_q=\sum_{n\ge1}2^{-n}\min(1,(q-h_n)_+)
\]
is positive exactly on this union. This proves the second condition.

Conversely choose an integer $M\ge\sup_Xu$. Set $v=e^{u-M}$, with
$e^{-\infty}=0$. For each rational $q<M$, choose continuous
$f_q:X\to[0,1]$ positive exactly on $\{u<q\}$, and put
\[
 r_q=e^{q-M},\qquad
 k_{q,j}=r_q+(1-r_q)(1-jf_q)_+\quad(j\ge1).
\]
Each $k_{q,j}$ is continuous, lies in $[r_q,1]$, and dominates $v$.
Indeed if $u\ge q$, then $f_q=0$ and $k_{q,j}=1$; if $u<q$, then
$v<r_q\le k_{q,j}$. For a point with $u<M$, choose $q>u$ arbitrarily
close to $u$ (arbitrarily negative if $u=-\infty$). Since $f_q>0$ there,
large $j$ gives $k_{q,j}=r_q$. At a point with $u=M$ all these functions
are one. Thus $\inf_{q,j}k_{q,j}=v$ everywhere. Enumerate the countable
pairs $(q,j)$ and let $K_n$ be the minimum of the first $n$ functions.
The functions $h_n=M+\log K_n$ are real continuous, decrease, and
converge everywhere to $u$, including at its poles. Finally,
$\{u<q\}$ is open under either condition; rational sublevels suffice
to prove upper semicontinuity.
\end{proof}

\section*{A concrete obstruction on the actual analytic carrier}
Let $g$ be the image of the joint unit Gauss point of $z_1,z_2$.
\begin{theorem}\label{nongdelta}
The singleton $\{g\}$ is not a $G_\delta$ subset of $X$.
Consequently the bounded upper semicontinuous function
\[
 u_g(y)=\begin{cases}1&y=g,\\0&y\ne g\end{cases}
\]
does not admit any sequence of real continuous functions decreasing
everywhere to it. In particular $u_g\notin\G$, irrespective of an
energy bound. This quantifies over the entire original class.
\end{theorem}
\begin{proof}
Use the ordinary torus open $V=\{|r_i|<1/4\}$. The deck element indexed
by $(m,n)\in\mathbb Z^2$ translates $r$ by $(m,n)$, since $2,3$ have
norm one and $|t|=e^{-1}$. Distinct translates of $V$ are disjoint.
The quotient topology therefore identifies $V$ homeomorphically with
its ordinary open image in $X$. This is only a local topological
assertion and supplies no coherent descent or nonalgebraization proof.

Consider any countable sequence of open neighbourhoods of $g$ and pull
them back to $V$. By the definition of the topology of the affine torus
analytification, each contains a basic neighbourhood specified by finitely
many conditions on norms of Laurent polynomials in
$K[z_1^{\pm1},z_2^{\pm1}]$. Collect the resulting countably many nonzero
Laurent polynomials $P_\ell$. Their norms at the joint unit Gauss point
are $M_\ell=\max|c_{\ell,I}|$. Choose a coefficient of that norm and
divide by it. All coefficients of the normalized polynomial lie in the
valuation ring, and its reduction is a nonzero Laurent polynomial over
the residue field $\mathbb C$.

Let $F\subset\mathbb C$ be the field over $\mathbb Q$ generated by all
coefficients of these countably many reductions. The field $F$ is
countable: it consists of finite rational expressions in a countable
set. The elements of $\mathbb C$ algebraic over a countable field form
a countable set, since there are countably many polynomials, each with
finitely many roots. Choose $\alpha\in\mathbb C$ transcendental over
$F$, and then $\beta\in\mathbb C$ transcendental over the countable
field $F(\alpha)$. This is possible because $\mathbb C$ is uncountable.
The nonzero constants $\alpha,\beta$ are algebraically independent over
$F$. Every reduced Laurent polynomial in the collected family is
nonzero at $(\alpha,\beta)$ (clear denominators by a monomial).

At the rigid torus point $p=(\alpha,\beta)$, its original coordinates
have norm one. Evaluation of each normalized polynomial has nonzero
residue, hence norm one. Thus $|P_\ell(p)|=M_\ell=|P_\ell(g)|$ for
every $\ell$. The point $p$ lies in $V$ and in every chosen basic
neighbourhood, hence in every member of the countable original sequence.
It differs from $g$: $z_1-\alpha$ has norm zero at $p$ and norm one at
$g$. Local injectivity preserves this distinction in the quotient.
No countable intersection of open neighbourhoods of $g$ can therefore
be its singleton.

Since $X$ is Hausdorff, $\{g\}$ is closed and $u_g$ is upper
semicontinuous. If $u_g$ had a continuous decreasing presentation,
Lemma \ref{barriers} would make $X\setminus\{g\}=\{u_g<1/2\}$ a
cozero set. Its complement would then be the zero set of a continuous
nonnegative $f$, equal to $\bigcap_{n\ge1}\{f<1/n\}$, a $G_\delta$.
This contradicts the first part. In particular no original-$\C$
presentation, which is a continuous presentation, exists.
\end{proof}

\section*{The exact whole-class obstacle profile}
For a bounded-above $u$ and real continuous $h$ with $h>u$ everywhere,
define, in the extended real numbers,
\[
 \Gamma_u(h)=\sup\{E(a):a\in\C,\ u\le a<h\text{ everywhere}\},
 \qquad
 \A(u)=\inf_{\substack{h\in C(X,\mathbb R)\\h>u}}\Gamma_u(h).
\]
The supremum of the empty set is $-\infty$. The obstacle family is
nonempty because $u$ is bounded above. The competitors are actual
members of the ENTIRE original class and use its cubic energy.
The profile $\A(u)$ is an obstacle-feasibility functional; it is not
a definition of, or a claim about, extended energy $E(u)$.

\begin{theorem}\label{recognition}
For every bounded-above $u:X\to[-\infty,\infty)$,
\[
 u\in\G\quad\Longleftrightarrow\quad
 \bigl[\{u<q\}\text{ is cozero for every rational }q\bigr]
 \ \text{and}\ \A(u)>-\infty.
\]
This is an exact recognition equivalence, not an assertion that the
conditions hold for a given raw limit. The forward direction requires
no monotonicity of the cubic energy; the reverse direction uses only
constant shifts and \eqref{shift}.
\end{theorem}
\begin{proof}
Let $a_n\downarrow u$ be an actual presentation with $E(a_n)\ge L$.
Lemma \ref{barriers} proves the cozero condition. Fix $h>u$ continuous.
The open sets $\{a_n<h\}$ increase and cover compact $X$. A finite
subcover yields one $n$ with $u\le a_n<h$ on all of $X$. Thus
$\Gamma_u(h)\ge L$ for every such obstacle and $\A(u)\ge L$.

Conversely obtain $h_n\downarrow u$ from Lemma \ref{barriers} and
choose a real $L<\A(u)$; this is possible also if $\A(u)=+\infty$.
Set $\epsilon_n=2^{-n-2}$ and
$s_n=\sum_{k=n}^{\infty}\epsilon_k=2^{-n-1}$.
For $H_1=h_1+\epsilon_1$ choose $b_1\in\C$ with
$u\le b_1<H_1$ and $E(b_1)>L$. This is possible since
$\Gamma_u(H_1)\ge\A(u)>L$.
Inductively the continuous obstacle
\[
 H_n=\min(h_n+\epsilon_n,b_{n-1}+\epsilon_{n-1})\quad(n\ge2)
\]
is strictly above $u$. Its minimum is only an ordinary continuous
obstacle; it is never asserted to be psh. Choose an actual $b_n\in\C$
with $u\le b_n<H_n$ and $E(b_n)>L$.
Then $b_n<h_n+\epsilon_n$ and
$b_{n+1}<b_n+\epsilon_n$. Define actual analytic weights
$a_n=b_n+s_n$. Since $s_n=\epsilon_n+s_{n+1}$,
\[
 a_{n+1}<a_n,\qquad
 u\le a_n<h_n+\epsilon_n+s_n,\qquad
 E(a_n)=E(b_n)+2s_n>L.
\]
The middle bounds imply $a_n\to u$ at every point, including where
$u=-\infty$. These are the required globally continuous decreasing
original-$\C$ weights with a uniform lower energy bound.
\end{proof}

\begin{corollary}\label{continuous}
For real continuous $u$, actual recognition in $\G$ is equivalent to
$u\in\C$, and also to $\A(u)>-\infty$.
\end{corollary}
\begin{proof}
Every continuous $u$ has the barrier $u+1/n$, so its cozero condition
holds. If $u\in\C$, the actual presentation $u+1/n$ has energies
$E(u)+2/n$, proving $u\in\G$. Conversely, for an actual $a_n\downarrow u$
and each $\delta>0$, the open sets $\{a_n<u+\delta\}$ increase and
cover compact $X$. Hence eventually $0\le a_n-u<\delta$ globally.
This proves uniform convergence and $u\in\C$ by its defining closure.
Theorem \ref{recognition} gives the remaining equivalence.
\end{proof}

\section*{All-support finite-dual moment control}
For a positive Radon measure $\mu$ of mass two retain
\[
 D(\mu)=\sup_{a\in\C}\left(E(a)-\int_Xa\,d\mu\right).
\]
This is FULL $D$. There is no support restriction.

\begin{theorem}\label{moment}
If $u$ satisfies the two recognition conditions and $D(\mu)<\infty$,
then $\A(u)$ is real, $\int_Xu\,d\mu$ is finite, and
\[
                   \int_Xu\,d\mu\ge\A(u)-D(\mu).
\]
In particular an eligible whole Dirac $2\delta_x$ cannot have an actual
finite-lower-energy presentation whose limit has a pole at $x$.
\end{theorem}
\begin{proof}
For every admissible $a<h$, the defining full-dual inequality and
positivity give $E(a)\le D(\mu)+\int_Xh\,d\mu$. Therefore
$\Gamma_u(h)\le D(\mu)+\int_Xh\,d\mu$, also when its competitor
set is empty. A constant obstacle bounds $\A(u)$ above by a finite
number, and the recognition hypothesis bounds it below.
Choose ordinary continuous $h_n\downarrow u$. The obstacles
$h_n+1/n$ are strictly above $u$, so
\[
 \A(u)\le D(\mu)+\int_Xh_n\,d\mu+2/n.
\]
Since $h_1$ is bounded, monotone convergence applied to $h_1-h_n$
gives $\int h_n\to\int u$ in $[-\infty,\infty)$. The real lower
bound in the last display rules out $-\infty$ and yields the claim.
For a whole Dirac the integral is $2u(x)$, proving the final statement.
\end{proof}

Constant translation preserves the criterion exactly:
$\Gamma_{u+c}(h+c)=\Gamma_u(h)+2c$ and
$\A(u+c)=\A(u)+2c$. These follow by the bijection $a\mapsto a+c$;
they keep raw mass two throughout.

\section*{Source applications and remaining original gaps}
The analytic structural applications above use Chambert--Loir--Ducros,
\emph{Formes differentielles reelles et courants sur les espaces de
Berkovich}, arXiv:1204.6277v2, header dated 25 July 2025; the retained
front matter also has an internal timestamp 29 July 2025, 0h06.
The exact applied assertions and
hypotheses, physical display ranges, and immediate-proof qualifications
are recorded in the separate JSON/TXT ledger. The whole immediate
units for Theorem 15.7.1; Definition 8.1.7 and Lemma 8.1.8; Lemmas
8.5.1, 8.5.3 and 8.5.5; construction 8.7.1 and Proposition 8.7.4;
Section 8.8.5; Proposition 7.2.8; and the cutoff/density chain 6.2.1--9
were physically read in bounded untruncated displays. Theorem 15.7.1
applies because the supplied surface is topologically separated and
pure-dimensional, with its actual empty analytic boundary. Its immediate
proof also establishes local differential separation on that interior,
as needed for 7.2.8. Section 8.5 assumes a boundaryless tropical space.
The local differences and smooth harmonic reference changes are exactly
those admitted by 8.8.5. No CLD 8.6.7 multiplier assertion is used.

The proofs of Lemma \ref{barriers}, Theorem \ref{recognition}, and
Theorem \ref{moment} are elementary and fully given here. They use
compactness, countability of the rationals, uniform convergence of
bounded continuous series, and ordinary monotone convergence, not a
classical psh recognition theorem or an external envelope theorem.
Theorem \ref{nongdelta} uses the defining seminorm topology of affine
analytification, the explicit deck translations, elementary reduction
of Laurent polynomials, and the proved countability facts about fields.
Its indicator is not asserted psh or an analytic equation potential;
the result is an actual whole-class presentation obstruction, not
nonexistence of a prescribed Monge--Ampere equation.
The deeper analytic references used inside the CLD units, including
Appendix A and the construction of the calibrated analytic carrier,
are not recursively audited. No incoming note or PASS verdict is a
correctness grant. The full inherited packet is retained separately;
its whole independent physical reading remains incomplete.

For the OLD target $x_{1/8}$, the isotropic Gauss point in
$(z_1-1,z_2-1)$ of radii $(e^{-1/8},e^{-1/8})$, this paper constructs
neither a target potential nor its strict-obstacle majorants. Even if
finite FULL $D(2\delta_{x_{1/8}})$ is independently established, that
inequality bounds $\Gamma_u$ from above; it does not establish its
uniform lower bound or feasibility. Recognition alone does not prove
that a recognized potential solves the prescribed equation. Variational
attainment, sufficient two-sided variations or orthogonality, comparison,
genuine nonalgebraization/descent, and the general existence-and-uniqueness
theorem are not supplied by this criterion. P16.7 and the unchanged
original residual remain OPEN. No other atom is claimed solved.
Candidate is false; independent mathematical correctness and compilation
remain unverified. ROOT's pending compiler requests are unchanged.
\end{document}
